% =============================================================================
% SVVVL Research Paper — Goal-Based Investing as Ergodic Control (Technical)
% Built from template/svvvl_paper_template.tex (shared house style).
% =============================================================================

\documentclass[11pt]{article}

% ----- Page geometry -------------------------------------------------------
\usepackage[
  letterpaper,
  top=0.9in,
  bottom=0.9in,
  left=0.85in,
  right=0.85in,
  headheight=14pt
]{geometry}

% ----- Encoding & language -------------------------------------------------
\usepackage[utf8]{inputenc}
\usepackage[T1]{fontenc}
\usepackage[english]{babel}



% ----- Typography ----------------------------------------------------------
\usepackage{XCharter}

% Section headings
\usepackage{titlesec}
\titleformat{\section}
  {\normalfont\Large\bfseries}{\thesection}{0.6em}{}
\titleformat{\subsection}
  {\normalfont\large\bfseries}{\thesubsection}{0.6em}{}
\titlespacing*{\section}{0pt}{1.4em}{0.6em}
\titlespacing*{\subsection}{0pt}{1.1em}{0.4em}

% Remove section numbering
\setcounter{secnumdepth}{0}

% ----- Spacing & paragraphs ------------------------------------------------
\usepackage{setspace}
\setstretch{1.15}
\setlength{\parskip}{0.6em}
\setlength{\parindent}{0pt}

% ----- Graphics & color ----------------------------------------------------
\usepackage{graphicx}
\usepackage{xcolor}
\definecolor{linkcolor}{HTML}{1F4E79}

% ----- Hyperlinks ----------------------------------------------------------
\usepackage[
  colorlinks=true,
  linkcolor=linkcolor,
  urlcolor=linkcolor,
  citecolor=linkcolor
]{hyperref}

% ----- Tables, figures, captions ------------------------------------------
\usepackage{booktabs}
\usepackage{array}
\usepackage{caption}
\captionsetup{font=small, labelfont=bf, skip=6pt}

% ----- Math ----------------------------------------------------------------
\usepackage{amsmath, amssymb}
\usepackage{amsthm}
\usepackage{mathtools}

% ----- Footnotes -----------------------------------------------------------
\usepackage[hang, flushmargin]{footmisc}
\setlength{\footnotesep}{0.8em}

% ----- Headers & footers ---------------------------------------------------
\usepackage{fancyhdr}
\pagestyle{fancy}
\fancyhf{}
\renewcommand{\headrulewidth}{0pt}
\fancyfoot[C]{\thepage}

\fancypagestyle{firstpage}{%
  \fancyhf{}%
  \renewcommand{\headrulewidth}{0pt}%
}


% ----- Logo lockup ---------------------------------------------------------
\newcommand{\WORDMARKFILE}{../assets/SVVVL_Research_Wordmark.pdf}
\newcommand{\WORDMARKWIDTH}{1.5in}
\newcommand{\LOCKUPGAP}{0.01in}

% The lockup itself: symbol above wordmark, both centered within a box
% the width of the symbol so the wordmark sits naturally beneath it.
\newcommand{\logolockup}{%
  \begin{minipage}{\WORDMARKWIDTH}
    \centering
    \includegraphics[width=\WORDMARKWIDTH]{\WORDMARKFILE}
  \end{minipage}%
}

% Title block command: \makecover{Title}{Subtitle}{Authors}{Date}
\newcommand{\makecover}[4]{%
  \thispagestyle{firstpage}%
  \noindent\logolockup\par
  \vspace{2.2em}
  {\fontsize{26pt}{30pt}\selectfont\bfseries
   \hyphenpenalty=10000\exhyphenpenalty=10000\sloppy
   #1\par}
  \vspace{0.6em}
  {\fontsize{15pt}{18pt}\selectfont\itshape\color{black!70}
   \hyphenpenalty=10000\exhyphenpenalty=10000\sloppy
   #2\par}
  \vspace{1.6em}
  {\normalsize\color{black!60} #3\par}
  \vspace{0.2em}
  {\normalsize\color{black!60} #4\par}
  \vspace{1.8em}
}

% ----- Theorem environments ------------------------------------------------
\newtheorem{proposition}{Proposition}
\newtheorem{remark}{Remark}
\newtheorem{conjecture}{Conjecture}
\newtheorem{theorem}{Theorem}
\newtheorem{lemma}{Lemma}
\newtheorem{corollary}{Corollary}
\newtheorem{definition}{Definition}

% ----- Custom commands -----------------------------------------------------
% General-purpose math operators reused across papers. Add paper-specific
% notation macros below as needed.
\newcommand{\E}{\mathbb{E}}
\newcommand{\Var}{\mathrm{Var}}
\newcommand{\Corr}{\mathrm{Corr}}
\newcommand{\dd}{\mathrm{d}}

% =============================================================================
% DOCUMENT
% =============================================================================
\begin{document}

% --- COVER -------------------------------------------------------------------
\makecover
  {Goal-Based Investing under Ergodic Dynamics}
  {Time-Average Growth, Geometric-Mean Targets, and a Fixed-Point Closure
   of the Constraint Identity}
  {J.~E. Whitehead~III\footnote{SVVVL, Towson, MD, United States.
   Contact: \href{mailto:james.whitehead@svvvl.com}{james.whitehead@svvvl.com}.
   See Disclosures at the end of this document.}}
  {\today}

% --- ABSTRACT ----------------------------------------------------------------
\begin{quote}
\itshape
We reformulate the goal-based portfolio problem under geometric
Brownian motion with the time-average growth rate of log wealth as
the objective and a geometric-mean terminal target as the constraint.
The reformulation produces an exact fixed-point closure of the
constraint identity: the contribution stream's effect on its own
discount base yields a future-value-of-annuity calculation in which
the time-average growth rate $R(l) = \mu_b + \mu_e l - \tfrac12
\sigma^2 l^2$ replaces the ensemble drift. The optimal allocation is
Kelly. The framework presents a one-parameter frontier of $(l, c)$
trade-offs that the investor traverses by preference between
volatility and contribution dollars. Model-predictive re-planning
from the realized state recovers the dynamic-programming optimum
exactly, because the constraint is linear in $\ln P$ and the
objective integrand is stationary.
\end{quote}


\newpage
% --- INTRODUCTION ----------------------------------------------------------
\section{Introduction}

The goal-based portfolio problem asks how an investor with current
wealth $P(t)$, terminal target $P_T$, horizon $T$, and a continuous
contribution stream $c(t)$ should allocate between a risky and a
risk-free asset. The standard formulation, following the goal-based
wealth management literature (Brunel, 2015; Das, Ostrov,
Radhakrishnan, and Srivastav, 2020), poses this as an expected-utility
maximization subject to an expected terminal wealth constraint
$\mathbb{E}[P(T)] = P_T$.

Both objects are ensemble averages. Under multiplicative wealth
dynamics, the ensemble drift $\mu_b + \mu_e l$ governs
$\mathbb{E}[P(t)]$, while the time-average growth rate of a single
path is $R(l) = \mu_b + \mu_e l - \tfrac12 \sigma^2 l^2$. The gap
between them is the volatility drag $\tfrac12 \sigma^2 l^2$, the
difference between the average outcome across paths and the typical
outcome along a path (Peters and Adamou, 2018; Peters, 2019). For an
investor who walks one path, the second number is the relevant one.

This note reformulates the problem in terms of objects that align
with the investor's experience. The objective becomes the expected
integral of $d\ln P$. The constraint becomes a geometric-mean target,
$\mathbb{E}[\ln P(T)] = \ln P_T$. Both are now phrased in $\ln P$,
which is the ergodic transformation for geometric Brownian motion:
the change of variables under which time and ensemble averages
of the increment process coincide.

The reformulation has four consequences.

First, the constraint identity admits an exact fixed-point closure.
Contributions feed back through the wealth base they build, and the
resulting expression for required contributions takes the form of a
standard future-value-of-annuity calculation, with the time-average
growth rate $R(l)$ replacing the ensemble drift.

Second, the optimal allocation is Kelly, $l^* = \mu_e/\sigma^2$, with
the savings rate as the residual control. The framework leaves the
allocation pinned at Kelly unless the investor's preferences over
volatility pull it away.

Third, the framework presents the investor a one-parameter frontier
of $(l, c)$ trade-offs at each rebalancing date. Deviations from
Kelly are admissible below $l^*$ as volatility-versus-savings
trade-offs and dominated above.

Fourth, the model-predictive re-planning structure that arises
naturally from this construction recovers the dynamic-programming
optimum exactly, because the constraint is linear in $\ln P$ and the
integrand of the objective is stationary.

The paper restricts attention to constant-parameter geometric
Brownian motion, where the ergodic transformation is the natural log
and the closed forms are elementary. Section~\ref{sec:setup} sets up
the dynamics and states the problem. Section~\ref{sec:identity}
derives the constraint identity and its fixed-point closure. Sections
\ref{sec:snapshot} and \ref{sec:lifecycle} present the two
user-facing frontiers, snapshot and lifecycle, and characterize
the Kelly minimum. Section~\ref{sec:feasibility} develops the
feasibility window under contribution and allocation bounds.
Section~\ref{sec:mpc} states the model-predictive algorithm and its
equivalence to dynamic programming. Section~\ref{sec:extensions}
discusses extensions. Derivations and proofs are in the appendix.

% --- SECTION ---------------------------------------------------------------
\section{Setup}
\label{sec:setup}

The portfolio value $P(t)$ evolves under
\begin{equation}
  \label{eq:dynamics}
  dP(t) \;=\; \bigl[P(t)(\mu_b + \mu_e l(t)) + c(t)\bigr]\,dt
  \;+\; P(t)\,\sigma\, l(t)\, dW_t,
  \qquad P(0) = P_0,
\end{equation}
where $W$ is a standard Brownian motion, $l(t)$ is the risky-asset
allocation, and $c(t) \geq 0$ is the contribution rate.

We assume throughout that the risk-free rate $\mu_b$, excess return
$\mu_e$, and volatility $\sigma$ are positive constants, and that the
horizon $T < \infty$ and target $P_T > 0$ are fixed.

Define the time-average log-growth rate at allocation $l$:
\begin{equation}
  \label{eq:R}
  R(l) \;\equiv\; \mu_b + \mu_e l - \tfrac12 \sigma^2 l^2.
\end{equation}
This is the deterministic drift of $\ln P$ in the absence of
contributions, and the rate that governs the typical path of $P$.

The investor's problem is
\begin{equation}
  \label{eq:problem}
  \max_{l(\cdot), c(\cdot)}\;
  \mathbb{E}\!\left[\int_0^T d\ln P(s)\right]
  \quad \text{subject to} \quad
  \mathbb{E}[\ln P(T)] = \ln P_T
\end{equation}
and the dynamics in~\eqref{eq:dynamics}. The objective and constraint
are both phrased in $\ln P$, which is the ergodic transformation for
GBM: $\ln P$ has stationary Gaussian increments, so time and ensemble
averages of $d\ln P$ coincide along any single realized path.

% --- SECTION ---------------------------------------------------------------
\section{The Constraint Identity and Its Closure}
\label{sec:identity}

\subsection{The constraint identity}

Applying It\^o's lemma to $\ln P$ under~\eqref{eq:dynamics}:
\begin{equation}
  \label{eq:dlnP}
  d\ln P \;=\; \left[R(l) + \frac{c}{P}\right]dt \;+\; \sigma l\, dW.
\end{equation}
Integrate from $t$ to $T$, take expectations conditional on
$\mathcal{F}_t$, and impose $\mathbb{E}[\ln P(T)] = \ln P_T$:
\begin{equation}
  \label{eq:identity}
  \ln(P_T/P(t)) \;=\; \int_t^T R(l(s))\,ds
  \;+\; \int_t^T \mathbb{E}\!\left[\frac{c(s)}{P(s)} \,\Big|\, \mathcal{F}_t\right]ds.
\end{equation}
The first term is the deterministic log-growth from the allocation
policy. The second is the conditional expected log-growth contributed
by savings. Equation~\eqref{eq:identity} is the working identity of
the framework.

\subsection{The fixed-point closure}

The second integrand in~\eqref{eq:identity} cannot be evaluated
directly because $P(s)$ in the denominator is the wealth process,
which itself depends on contributions made between $t$ and $s$.
Define the cumulative contribution-induced log-growth
\begin{equation}
  \label{eq:phi-def}
  \phi(s) \;\equiv\; \int_t^s \overline{c(u)/P(u)}\,du,
\end{equation}
where $\overline{c/P}$ denotes the time-average expectation along the
typical path. The self-reference in~\eqref{eq:phi-def} closes into a
separable ODE for $\phi$, with solution
\begin{equation}
  \label{eq:phi}
  \phi(s) \;=\; \ln\!\left[1 + \frac{c}{P(t)}\,A(R(l), s-t)\right],
  \qquad
  A(R, \tau) \;\equiv\; \frac{1 - e^{-R\tau}}{R},
\end{equation}
when $l$ and $c$ are held constant on $[t, T]$. The function $A(R,
\tau)$ is the standard continuous-time annuity factor at rate $R$
over horizon $\tau$. The derivation is in Appendix~\ref{app:phi}.

Substituting $\phi(T)$ into~\eqref{eq:identity}, exponentiating, and
rearranging produces the working form of the constraint, which we
call the \emph{dollar-balance equation}:
\begin{equation}
  \label{eq:balance}
  \boxed{
  P(t)\, e^{R(l)(T-t)}
  \;+\; c\cdot \frac{e^{R(l)(T-t)} - 1}{R(l)}
  \;=\; P_T.
  }
\end{equation}
This is the standard future-value-of-annuity identity from finance:
current wealth $P(t)$ compounded forward at rate $R(l)$, plus a
continuous-flow annuity $c$ compounded forward at the same rate,
equals the terminal target $P_T$. The rate $R(l)$ is the time-average
growth rate of log wealth, not the ensemble drift $\mu_b + \mu_e l$,
so the volatility drag is built into the compounding by construction.
The derivation is in Appendix~\ref{app:balance}.

The textbook future-value-of-annuity formula uses $\mu_b + \mu_e l$
as the rate. The replacement of this rate by $R(l)$ is the working
content of the reformulation. Numerically, for $\sigma = 0.18$ and
$l = 1$, the difference is $\tfrac12 \sigma^2 l^2 = 1.62\%$ per year,
which compounds to a thirty-percent gap over a twenty-year horizon.

% --- SECTION ---------------------------------------------------------------
\section{The Snapshot Frontier}
\label{sec:snapshot}

Solving~\eqref{eq:balance} for $c$ in terms of $l$ gives the
snapshot frontier:
\begin{equation}
  \label{eq:modeA}
  \boxed{
  c(l) \;=\;
  \frac{P_T \;-\; P(t)\,e^{R(l)(T-t)}}{\bigl(e^{R(l)(T-t)} - 1\bigr) / R(l)}.
  }
\end{equation}
The frontier is the set of $(l, c)$ pairs that, held constant from
$t$ to $T$, satisfy the geometric-mean terminal constraint.

Whenever the goal exceeds the Kelly-compounded current wealth (that
is, $P_T > P(t)\, e^{R(l^*)(T-t)}$ with $l^* = \mu_e/\sigma^2$),
$c(l)$ is strictly convex on $[l_{\min}, l_{\max}]$ with global
minimum at the Kelly allocation $l^*$. The Kelly allocation maximizes
$R(l)$ over $\mathbb{R}$, which minimizes the compounded shortfall in
the numerator and the annuity factor in the denominator; the former
effect dominates. The derivation is in Appendix~\ref{app:convex}.

For $l > l^*$, both $c(l)$ and $\mathrm{Var}(P(T))$ are increasing in
$l$. Allocations above Kelly are therefore dominated under the joint
criterion of minimizing required contributions and minimizing
portfolio variance. The investor's admissible choice set is $l \in
[l_{\min}, l^*]$. The trade-off along this segment is between
allocation volatility and contribution level: lower $l$ reduces
portfolio variance at the cost of higher $c$.

% --- SECTION ---------------------------------------------------------------
\section{The Lifecycle Frontier}
\label{sec:lifecycle}

A more realistic specification lets contributions grow with income.
Set
\begin{equation}
  c(s) \;=\; c_0\, e^{g s},
\end{equation}
with $c_0 > 0$ and $g \in \mathbb{R}$ deterministic. The same
fixed-point argument applied to this deterministic contribution
schedule produces the lifecycle constraint:
\begin{equation}
  \label{eq:modeB}
  \boxed{
  P(t)\, e^{R(l)(T-t)}
  \;+\; c_0\, e^{gt}\, e^{R(l)(T-t)}\cdot \frac{e^{(g - R(l))(T-t)} - 1}{g - R(l)}
  \;=\; P_T.
  }
\end{equation}
The case $g = R(l)$ is handled by continuity:
$(e^{(g-R(l))(T-t)} - 1)/(g - R(l)) \to T - t$. The derivation is in
Appendix~\ref{app:lifecycle}.

Solving for $c_0$ at fixed $(g, l)$:
\begin{equation}
  \label{eq:c0}
  c_0 \;=\;
  \frac{\bigl[P_T\, e^{-R(l)(T-t)} - P(t)\bigr]\,
        e^{-gt}\,(g - R(l))}
       {e^{(g - R(l))(T-t)} - 1}.
\end{equation}
The investor specifies the wage-growth assumption $g$ and the
allocation $l$, and reads off the required initial contribution.

At fixed $g \neq R(l^*)$, $c_0$ is strictly convex in $l$ with
minimum at $l^* = \mu_e/\sigma^2$ on the interior of the feasibility
region. The same dominance argument as for the snapshot frontier
restricts the investor's admissible choice set to $l \in [l_{\min},
l^*]$. The derivation is in Appendix~\ref{app:lifecycle-convex}.

% --- SECTION ---------------------------------------------------------------
\section{The Feasibility Window}
\label{sec:feasibility}

The investor's allocation is bounded by risk-tolerance limits
$l \in [l_{\min}, l_{\max}]$ and her contribution by dollar capacity
$c \in [\text{LB}, \text{UB}]$. These bounds imply, at any state
$(P(t), t)$, a window of reachable terminal wealth under the
hold-to-expiry assumption.

Define $l_{\text{best}} = \min(l_{\max}, \mu_e/\sigma^2)$, and let
$l_{\text{worst}} \in \{l_{\min}, l_{\max}\}$ denote the endpoint
minimizing $R(l)$ over $[l_{\min}, l_{\max}]$. With constant $l$ and
constant $c$ held to expiry, the set of reachable terminal wealth
values is $P_T \in [P(T)_{\min}, P(T)_{\max}]$, where
\begin{align}
  \label{eq:Pmax}
  P(T)_{\max} &= P(t)\, e^{R(l_{\text{best}})(T-t)}
  \;+\; \mathrm{UB} \cdot \frac{e^{R(l_{\text{best}})(T-t)} - 1}{R(l_{\text{best}})}, \\
  \label{eq:Pmin}
  P(T)_{\min} &= P(t)\, e^{R(l_{\text{worst}})(T-t)}
  \;+\; \mathrm{LB} \cdot \frac{e^{R(l_{\text{worst}})(T-t)} - 1}{R(l_{\text{worst}})}.
\end{align}

The window expands when $P(t)$ rises and contracts when it falls,
each rebalancing period. A goal $P_T > P(T)_{\max}$ is infeasible
from the current state and signals a required revision of $P_T$, $T$,
or the bounds.

% --- SECTION ---------------------------------------------------------------
\section{Model-Predictive Algorithm and Bellman Equivalence}
\label{sec:mpc}

\subsection{Algorithm}

At each rebalancing time $t \in [0, T)$:

\begin{enumerate}
  \item Observe the realized portfolio value $P(t)$.
  \item Compute the feasibility window $[P(T)_{\min}, P(T)_{\max}]$ via~\eqref{eq:Pmax} and~\eqref{eq:Pmin}.
  \item If $P_T \notin [P(T)_{\min}, P(T)_{\max}]$, prompt for
    revision of $P_T$, $T$, or bounds.
  \item Otherwise, compute the chosen frontier:~\eqref{eq:modeA}
    (snapshot) or~\eqref{eq:c0} (lifecycle).
  \item The investor selects $(l, c)$ on the frontier subject to
    bounds.
  \item Hold the policy over $[t, t + \Delta]$. Return to step 1.
\end{enumerate}

Every step is closed-form. The control rule itself involves no
approximation of $1/P(s)$.

\subsection{Bellman equivalence}

Let $J^*(P(t), t)$ denote the optimal value of~\eqref{eq:problem}
under the geometric-mean constraint from initial state $(P(t), t)$.
The policy obtained by re-solving~\eqref{eq:identity} at each $t$ from
the realized $P(t)$ coincides almost surely with the optimal feedback
control of the dynamic-programming problem with value function $J^*$.
The argument turns on two properties of the reformulated problem:
the integrand of the objective is stationary in $(l, c/P)$, and the
constraint is linear in $\ln P$. Together, these properties make the
re-solved problem from $(P(t), t)$ the dynamic-programming subproblem
of the original. The derivation is in Appendix~\ref{app:bellman}.

Under the arithmetic-mean constraint $\mathbb{E}[P(T)] = P_T$, the
analogous equivalence fails: the variational re-solution is on
$\mathbb{E}[P]$ but the HJB optimum is on $\mathbb{E}[\ln P]$, and
the Jensen gap is exactly the volatility drag absorbed by $R(l)$ in
the geometric-mean formulation.

% --- SECTION ---------------------------------------------------------------
\section{Extensions}
\label{sec:extensions}

\textbf{Multiple risky assets.} For an $n$-vector of risky assets with
excess return $\mu_e \in \mathbb{R}^n$ and covariance $\Sigma \in
\mathbb{R}^{n \times n}$, the Kelly allocation generalizes to $l^* =
\Sigma^{-1}\mu_e$, with $R(l) = \mu_b + \mu_e^\top l - \tfrac12
l^\top \Sigma l$. The dollar-balance form~\eqref{eq:balance} and the
fixed-point closure are unchanged.

\textbf{Time-varying parameters.} Piecewise-constant $(\mu_b, \mu_e,
\sigma)$ extend the framework by recomputing the constraint identity
each period. The closed forms hold within each segment.

\textbf{Merton recovery.} Setting $P_T = P_0 \exp(r T)$ for
$r = \mu_b + \mu_e^2/(2\sigma^2)$ (the Kelly-compounded current
wealth) gives $c = 0$ at the Kelly allocation; the framework reduces
to the unconstrained log-optimal portfolio.

\textbf{Anomalous diffusions.} The natural log is the ergodic
transformation specific to geometric Brownian motion. For
subdiffusive, superdiffusive, or heavy-tailed processes, the
corresponding generalized log replaces $\ln$ throughout, and the
fixed-point closure carries over with the appropriate annuity factor.
A companion note develops this generalization.

% --- CONCLUSION ------------------------------------------------------------
\section{Conclusion}
\label{sec:conclusion}

The reformulation of the goal-based portfolio problem around the
time-average growth rate of log wealth and a geometric-mean terminal
target produces a framework with three operational properties: an
exact fixed-point closure of the constraint identity, a closed-form
$(l, c)$ frontier with a Kelly minimum, and a model-predictive
re-planning loop that recovers the dynamic-programming optimum
exactly. The framework respects the experience of an investor on a
single realized path rather than the expectation across a population
of paths she does not traverse. The implementation is elementary at
each rebalancing date, and the same construction extends to multiple
assets, time-varying parameters, and anomalous diffusions through the
appropriate generalized log.

% --- TECHNICAL APPENDIX ----------------------------------------------------
\newpage
\section{Technical Appendix}
\label{app:proofs}

\subsection{Derivation of the fixed-point closure}
\label{app:phi}

Let $m(s) = \mathbb{E}[\ln P(s) \mid \mathcal{F}_t]$. Taking
conditional expectations of~\eqref{eq:dlnP}:
\begin{equation}
  m'(s) \;=\; R(l) \;+\; \overline{c/P(s)}.
\end{equation}
Define $\overline{1/P(s)} \equiv 1/\exp(m(s))$, the ergodic
expectation of $1/P(s)$ under the time-average framing. Since the
contribution $c(s) = c$ is deterministic,
\begin{equation}
  \overline{c/P(s)} \;=\; c \cdot \overline{1/P(s)} \;=\; \frac{c}{e^{m(s)}}.
\end{equation}
Therefore
\begin{equation}
  m'(s) \;=\; R(l) \;+\; c\,e^{-m(s)},
\end{equation}
with $m(t) = \ln P(t)$. Substituting $\phi(s) = m(s) - \ln P(t) -
R(l)(s - t)$ so that $\phi(t) = 0$ and $m(s) = \ln P(t) + R(l)(s - t)
+ \phi(s)$:
\begin{equation}
  \phi'(s) \;=\; c\,e^{-\ln P(t) - R(l)(s - t) - \phi(s)}
  \;=\; \frac{c}{P(t)}\, e^{-R(l)(s - t)}\, e^{-\phi(s)}.
\end{equation}
This is a separable ODE. Multiply both sides by $e^{\phi(s)}$:
\begin{equation}
  e^{\phi(s)}\phi'(s) \;=\; \frac{c}{P(t)}\, e^{-R(l)(s - t)}.
\end{equation}
The left side equals $\frac{d}{ds}\bigl[e^{\phi(s)}\bigr]$. Integrating
from $t$ to $s$ with $e^{\phi(t)} = 1$:
\begin{equation}
  e^{\phi(s)} - 1 \;=\; \frac{c}{P(t)}\int_t^s e^{-R(l)(u - t)}\,du
  \;=\; \frac{c}{P(t)}\, A(R(l), s - t).
\end{equation}
Solving for $\phi(s)$ gives~\eqref{eq:phi}.

\subsection{Derivation of the dollar-balance equation}
\label{app:balance}

Substituting $\phi(T)$ from~\eqref{eq:phi} into the constraint
identity~\eqref{eq:identity} with constant $l$:
\begin{equation}
  \ln(P_T/P(t)) \;=\; R(l)(T - t) \;+\; \ln\!\left[1 + \frac{c}{P(t)}\,A(R(l), T - t)\right].
\end{equation}
Exponentiating:
\begin{equation}
  \frac{P_T}{P(t)}\, e^{-R(l)(T - t)} \;=\; 1 + \frac{c}{P(t)}\, A(R(l), T - t).
\end{equation}
Multiplying both sides by $P(t)\,e^{R(l)(T - t)}$:
\begin{equation}
  P_T \;=\; P(t)\, e^{R(l)(T - t)} \;+\; c\, e^{R(l)(T - t)}\, A(R(l), T - t).
\end{equation}
Using $e^{R\tau} A(R, \tau) = (e^{R\tau} - 1)/R$ gives~\eqref{eq:balance}.


\subsection{Convexity of the snapshot frontier}
\label{app:convex}

Write $\tau = T - t$ and $S(l) = (e^{R(l)\tau} - 1)/R(l)$ for the
future-value annuity factor. From~\eqref{eq:modeA},
\begin{equation}
  c(l) \;=\; \frac{P_T - P(t)\,e^{R(l)\tau}}{S(l)}.
\end{equation}
Differentiating with respect to $l$ and using $R'(l) = \mu_e -
\sigma^2 l$:
\begin{equation}
  c'(l) \;=\; -\frac{R'(l)}{S(l)^2}
  \left[P(t)\,e^{R(l)\tau}\,\tau\,S(l)
        + \bigl(P_T - P(t)\,e^{R(l)\tau}\bigr)\,S'(l) / R'(l) \cdot R'(l)\right].
\end{equation}
A cleaner route is to differentiate the implicit form. From the
dollar balance~\eqref{eq:balance},
\begin{equation}
  P(t)\, e^{R(l)\tau} + c(l)\, S(l) = P_T.
\end{equation}
Differentiating:
\begin{equation}
  P(t)\,\tau\, e^{R(l)\tau}\, R'(l) + c'(l)\, S(l) + c(l)\, S'(l) = 0,
\end{equation}
so
\begin{equation}
  c'(l) \;=\; -\frac{P(t)\,\tau\, e^{R(l)\tau}\, R'(l) + c(l)\, S'(l)}{S(l)}.
\end{equation}
At $l^* = \mu_e/\sigma^2$, $R'(l^*) = 0$ and $S'(l^*) = 0$ (the latter
because $S$ depends on $l$ only through $R(l)$, and $R'(l^*) = 0$).
Therefore $c'(l^*) = 0$. The second-order condition reduces, after
similar substitution and using $R''(l) = -\sigma^2$, to
\begin{equation}
  c''(l^*) \;=\; -\frac{P(t)\,\tau\, e^{R(l^*)\tau}\, R''(l^*)}{S(l^*)}
  \;=\; \frac{P(t)\,\tau\, e^{R(l^*)\tau}\, \sigma^2}{S(l^*)} \;>\; 0
\end{equation}
under the stated condition $P_T > P(t)\,e^{R(l^*)\tau}$ (so $c(l^*) >
0$ and $S(l^*) > 0$). Strict convexity follows by continuity and the
absence of other critical points on $[l_{\min}, l_{\max}]$ since
$R'(l)$ vanishes only at $l^*$.

\subsection{Derivation of the lifecycle constraint}
\label{app:lifecycle}

The fixed-point argument of Appendix~\ref{app:phi} extends to
deterministic $c(s) = c_0 e^{gs}$. The ODE for $\phi$ becomes
\begin{equation}
  \phi'(s) \;=\; \frac{c_0 e^{gs}}{P(t)}\, e^{-R(l)(s - t)}\, e^{-\phi(s)}.
\end{equation}
Multiplying by $e^{\phi(s)}$ and integrating:
\begin{equation}
  e^{\phi(s)} - 1 \;=\; \frac{c_0}{P(t)}\int_t^s e^{(g - R(l))(u - t) + R(l)(0)}\, du,
\end{equation}
which simplifies to
\begin{equation}
  e^{\phi(s)} - 1 \;=\; \frac{c_0\, e^{gt}}{P(t)}\cdot \frac{e^{(g - R(l))(s - t)} - 1}{g - R(l)}.
\end{equation}
At $s = T$, substituting $\phi(T)$ into~\eqref{eq:identity},
exponentiating, and multiplying through by $P(t)\,e^{R(l)(T - t)}$
yields~\eqref{eq:modeB}.

\subsection{Convexity of the lifecycle frontier}
\label{app:lifecycle-convex}

Differentiating~\eqref{eq:c0} implicitly through the dollar balance,
the same structure as Appendix~\ref{app:convex} applies: $c_0$ is
convex in $l$ because the dominant $l$-dependence enters through
$R(l)$, and $R$ is concave with maximum at Kelly. The
implicit-function argument carries over with $R'(l^*) = 0$ giving
$c_0'(l^*) = 0$ and $R''(l^*) = -\sigma^2$ giving $c_0''(l^*) > 0$
on the interior of feasibility.

\subsection{Bellman equivalence under the geometric-mean constraint}
\label{app:bellman}

The dynamic-programming subproblem from state $(P(t), t)$ is
\begin{equation}
  J^*(P(t), t) \;=\; \sup_{l, c}\, \mathbb{E}\!\left[\int_t^T d\ln P(s) \,\Big|\, \mathcal{F}_t\right]
  \quad \text{s.t.} \quad \mathbb{E}[\ln P(T) \mid \mathcal{F}_t] = \ln P_T.
\end{equation}
Two observations close the proof.

First, the objective integrand $d\ln P/dt = R(l) + c/P$ depends on
the controls $(l, c)$ and the state $P$ only through stationary
functions: $R(l)$ depends on $l$ alone, and $c/P$ on the
contemporaneous control and state. There is no explicit time
dependence and no path dependence through $\ln P$ history.

Second, the constraint is linear in $\ln P(T)$, with $\ln P(t)$
entering only as a shift: $\mathbb{E}[\ln P(T) \mid \mathcal{F}_t] -
\ln P(t) = \ln(P_T/P(t))$, which depends on the realized state $P(t)$
but not on the path leading to it.

Together, these properties imply that the constraint identity
$\eqref{eq:identity}$ rewritten from $(P(t), t)$ has the same
functional form as from $(P_0, 0)$, with $P(t)$ in place of $P_0$ and
$T - t$ as the remaining horizon. The re-solved policy is therefore
the optimal continuation by the principle of optimality (Bellman,
1957).

% --- REFERENCES ------------------------------------------------------------
\newpage
\section{Further Reading and References}

\begin{itemize}\setlength{\itemsep}{0.3em}
  \item Bellman, R.~E. (1957). \textit{Dynamic Programming}. Princeton
        University Press, Princeton.
  \item Brunel, J.~L.~P. (2015). \textit{Goals-Based Wealth Management:
        An Integrated and Practical Approach to Changing the Structure of
        Wealth Advisory Practices}. Wiley Finance, Hoboken, NJ.
  \item Cover, T.~M. (1991). ``Universal portfolios.''
        \textit{Mathematical Finance}, 1(1), 1--29.
  \item Das, S.~R., Ostrov, D., Radhakrishnan, A., and Srivastav, D.
        (2020). ``Dynamic portfolio allocation in goals-based wealth
        management.'' \textit{Computational Management Science}, 17(4),
        613--640.
  \item Kelly, J.~L. (1956). ``A new interpretation of information
        rate.'' \textit{Bell System Technical Journal}, 35(4), 917--926.
  \item Merton, R.~C. (1971). ``Optimum consumption and portfolio rules
        in a continuous-time model.'' \textit{Journal of Economic
        Theory}, 3(4), 373--413.
  \item Peters, O. (2019). ``The ergodicity problem in economics.''
        \textit{Nature Physics}, 15(12), 1216--1221.
  \item Peters, O., and Adamou, A. (2018). \textit{Ergodicity
        Economics}. Lecture notes, London Mathematical Laboratory.
  \item Peters, O., and Gell-Mann, M. (2016). ``Evaluating gambles
        using dynamics.'' \textit{Chaos}, 26(2), 023103.
\end{itemize}

\newpage
% --- DISCLOSURES -------------------------------------------------------------
\section{Disclosures}

The views expressed in this document are those of the author and
SVVVL alone and do not constitute investment, legal, tax, or
accounting advice. This document is provided for informational and
research purposes only and is not an offer or solicitation to buy or
sell any security or financial product.

The framework presented is theoretical and relies on idealized
assumptions that do not hold in practice. The results
should not be implemented in a real portfolio without additional
constraints and considerations not addressed in this paper.

While reasonable care has been taken in the preparation of this
document, no representation or warranty, express or implied, is made
as to its accuracy or completeness.

Reading this document does not create an advisory, fiduciary, or
client relationship of any kind between the reader and SVVVL or the
author.

\end{document}
